\section{Thảo luận và hạn chế}

Mô hình hiện tại đạt $0,8697$ độ chính xác trên tập kiểm thử với một kiến trúc nhỏ gọn $1,38$ triệu tham số, một con số tương đương các baseline GRU đã công bố trên IMDB nhưng vẫn còn khoảng cách đáng kể so với các kiến trúc Transformer hiện đại đạt mức trên $0,95$. Khoảng cách này xuất phát từ nhiều nguồn:

\begin{enumerate}
    \item \textbf{Giới hạn độ dài văn bản:} Ngưỡng cắt $256$ token loại bỏ một phần đáng kể thông tin của những bài đánh giá rất dài, đặc biệt là những bài có kết luận được đặt ở đoạn cuối. Mở rộng ngưỡng lên $512$ hoặc $1024$ sẽ làm chi phí tính toán tăng tuyến tính trong trường hợp GRU và tăng theo bậc hai trong trường hợp Transformer, đặt ra một bài toán đánh đổi (\textit{trade-off}) không tầm thường về mặt tài nguyên.
    \item \textbf{Giới hạn quy mô từ vựng:} Từ vựng $10\,000$ token loại bỏ phần lớn các thực thể có tên (tên phim, tên diễn viên), trong đó một số mang tín hiệu cảm xúc cực mạnh — ví dụ như tên của một series \textit{cult-classic} — và việc thay thế chúng bằng \texttt{<unk>} là một mất mát thông tin khó định lượng.
    \item \textbf{Tính ngẫu nhiên trong thực nghiệm:} Kết quả báo cáo hiện tại chỉ dựa trên một hạt giống ngẫu nhiên (\textit{random seed}) duy nhất. Một nghiên cứu khoa học chặt chẽ hơn cần báo cáo giá trị trung bình và độ lệch chuẩn qua ít nhất ba tới năm hạt giống để loại trừ các kết luận sai lệch do dao động ngẫu nhiên của một lần chạy.
    \item \textbf{Không gian tìm kiếm siêu tham số hẹp:} Nghiên cứu siêu tham số trong báo cáo này mới chỉ hạn chế ở ba điểm rời rạc trong không gian hai chiều $(\text{learning rate}, \text{batch size})$. Một quy trình chuyên sâu hơn nên áp dụng tìm kiếm Bayesian (\textit{Bayesian Optimization}) hoặc \textit{Hyperband} để khám phá một cách có hệ thống không gian này, kết hợp với các yếu tố cốt lõi khác như chiều \textit{embedding} và số lớp GRU.
\end{enumerate}

Hướng phát triển trực tiếp nhất là triển khai hai biến thể kiến trúc đã đề xuất ở mục 8, tiến hành so sánh đối chiếu sòng phẳng (\textit{apples-to-apples}) với baseline trên cùng một quy trình huấn luyện. Một hướng bổ sung có giá trị thực tiễn cao là khởi tạo lớp \textit{embedding} bằng các vectơ tiền huấn luyện (\textit{pretrained word embeddings}) như GloVe hoặc FastText. Việc này giúp mô hình tận dụng được kiến thức từ vựng học được trên các kho ngữ liệu (\textit{corpus}) lớn hơn nhiều lần so với IMDB, qua đó giảm sự phụ thuộc vào dữ liệu trong miền (\textit{in-domain}) và có thể cải thiện đáng kể cả tốc độ hội tụ lẫn chất lượng phân loại cuối cùng.